\reading{IM-9}{Risk, Regulation and Organizational Structure}{DEFER}{0}

\srcnote{David P. Stowell and Paul Stowell, \textit{Investment Banks, Hedge
Funds, and Private Equity}, 4th Edition (Elsevier, 2023), Chapter 14. No GARP
source book covers this reading; see the note at the head of IM-8.}

\begin{imkeybox}[title={The four objectives in this reading}]
\begin{enumerate}[label=\textbf{\alph*.}]
\item Describe different risks faced by investors in hedge funds.
\item Explain how the activities of hedge funds and exposure to hedge funds can increase systemic risk.
\item Summarize different approaches taken by regulators to mitigate risks related to hedge funds.
\item Describe a typical organizational structure for a hedge fund and explain motivations for adopting this structure.
\end{enumerate}
\end{imkeybox}

\begin{imtrapbox}[breakable=false,title={An objective with no label on it}]
The source notes label objectives a, b and d, and omit c entirely. The material
for c is present --- it runs across the regulation pages and the regulatory
objectives section --- but it carries no objective heading, so it reads as a gap
and is not one. It is attached to IM-9 c below.
\end{imtrapbox}

% ------------------------------------------------------------------
\lo{IM-9 a}{Describe different risks faced by investors in hedge funds.}

Hedge fund risks can be classified into \term{portfolio-level} and
\term{investment-level} risks.

\begin{itemize}
\item \term{Portfolio-level:} benchmarking issues, survivorship bias, and
\term{unrelated business taxable income (UBTI)}, which may arise when
debt-financed income becomes taxable for otherwise tax-exempt investors.
\item \term{Investment-level:} counterparty risk, conflicts of interest,
\term{key person risk}, strategy failure, \term{style drift}, regulatory changes,
and correlation spikes.
\end{itemize}

\subsection{Leverage}

Hedge funds frequently amplify outcomes through leverage, which can take the form
of \term{on-balance-sheet} leverage (direct borrowing) or
\term{off-balance-sheet} leverage, for example derivatives.

\figwrap{%
\begin{tikzpicture}
\begin{axis}[frmaxis,
  width=11.4cm, height=5.2cm,
  ybar, bar width=17pt,
  symbolic x coords={Relative value, Multistrategy, Macro, Equity hedge,
                     Event driven, Credit, Managed futures},
  xtick=data,
  x tick label style={font=\scriptsize\sffamily, rotate=28, anchor=north east},
  enlarge x limits=0.09,
  ymin=0, ymax=3.5, ytick={0,1,2,3},
  ylabel={average leverage ($\times$ NAV)},
]
\addplot[draw=FRMNavy,fill=PaleNavy,area legend] coordinates
  {(Relative value,3.0) (Multistrategy,2.2) (Macro,2.1) (Equity hedge,1.6)
   (Event driven,1.5) (Credit,1.5) (Managed futures,1.4)};
\node[font=\scriptsize\sffamily,color=black!70,anchor=south] at (axis cs:Relative value,3.05) {$3.0$};
\node[font=\scriptsize\sffamily,color=black!70,anchor=south] at (axis cs:Multistrategy,2.25) {$2.2$};
\node[font=\scriptsize\sffamily,color=black!70,anchor=south] at (axis cs:Macro,2.15) {$2.1$};
\node[font=\scriptsize\sffamily,color=black!70,anchor=south] at (axis cs:Equity hedge,1.65) {$1.6$};
\node[font=\scriptsize\sffamily,color=black!70,anchor=south] at (axis cs:Event driven,1.55) {$1.5$};
\node[font=\scriptsize\sffamily,color=black!70,anchor=south] at (axis cs:Credit,1.55) {$1.5$};
\node[font=\scriptsize\sffamily,color=black!70,anchor=south] at (axis cs:Managed futures,1.45) {$1.4$};
\end{axis}
\end{tikzpicture}}%
{Hedge fund leverage by strategy, from 2020 findings of the U.S. Department of
the Treasury's Office of Financial Research (source notes, Figure 91.2). Redrawn
rather than reproduced: the original is a scan with a sloping baseline that makes
the bars hard to compare, while every value is printed on it. \term{Relative
value carries twice the leverage of managed futures} --- and relative value is
precisely the family of strategies from IM-8 b whose individual trades are
described as ``low risk per trade, small but consistent returns.'' The leverage
is how those small returns are made worth having, and it is also why they are not
low-risk at the fund level.}

\subsection{Short selling, transparency and risk tolerance}

\begin{itemize}
\item \term{Short selling.} While long positions have limited downside, short
positions have \term{theoretically unlimited losses} because prices can rise
indefinitely. A \term{short squeeze} occurs when a heavily shorted stock rises
quickly, forcing short sellers to buy shares to cover, which pushes prices up
further --- as in the GameStop episode.
\item \term{Transparency.} Hedge funds often keep strategies secret to preserve
arbitrage opportunities. That limited transparency makes it harder to monitor
risks and ensure strategy alignment, and \term{exit gates and liquidity
restrictions} can prevent investors from leaving when they are concerned about
strategy changes.
\item \term{Risk tolerance.} Managers generally have higher risk tolerance than
mutual fund managers due to \term{compensation incentives}, and invest in a wider
range of complex strategies using derivatives, which can help manage risk but also
increase the complexity of assessing it.
\end{itemize}

% ------------------------------------------------------------------
\lo{IM-9 b}{Explain how the activities of hedge funds and exposure to hedge funds can increase systemic risk.}

\subsection{Systemic risk from hedge fund activities}

\begin{enumerate}
\item \term{Contagion through forced liquidation.} If a fund suffers large losses
it may unwind positions quickly and sell at \term{fire-sale prices}, pushing asset
prices down and spreading losses to other institutions holding similar assets.
\item \term{Losses to financial counterparties.} Funds borrow from banks and trade
derivatives with financial institutions; if a fund fails, banks and counterparties
may face significant losses.
\end{enumerate}

\begin{imkeybox}[title={Two failures, two different verdicts}]
\term{LTCM (1998)} highlighted the dangers of excessive leverage and liquidity
risk, and showed that assets which normally move independently \term{can become
highly correlated during periods of extreme market stress}.

\smallskip
\term{Amaranth Advisors (2006)} raised the same concerns, but the \emph{Bank of
England concluded that major financial intermediaries contribute more to systemic
risk than hedge funds}, and that hedge funds can sometimes help distribute risk in
the financial system. Critics still argue that high leverage and complex
strategies amplify systemic risk.
\end{imkeybox}

\subsection{Systemic risk from bank exposure to hedge funds}

\renewcommand{\arraystretch}{1.2}
\noindent\begin{tabularx}{\textwidth}{@{}>{\raggedright\arraybackslash}p{2.6cm}X@{}}
\toprule
\textbf{\sffamily\small Risk to the bank} & \textbf{\sffamily\small How it arises} \\
\midrule
Revenue risk & Banks earn fees from trading, clearing, custody, securities lending, financing and reporting. If multiple hedge funds fail, banks lose a significant source of revenue. \\
Credit risk & Banks provide margin loans secured by collateral. If the collateral value falls rapidly, the loan becomes \term{undercollateralized}. \\
Counterparty risk & Banks enter derivative contracts --- CDSs and other OTC derivatives. If a fund defaults, the bank faces losses on those contracts. \\
\bottomrule
\end{tabularx}
\renewcommand{\arraystretch}{1}
\figcap{Three distinct exposures, often collapsed into one in a question. Revenue
risk does not require any fund to default --- only for the business to disappear.
Credit risk requires collateral to fall faster than the bank can call for more.
Counterparty risk requires an actual default.}

\begin{imkeybox}[title={Archegos, 2021}]
The collapse of Archegos Capital highlighted the risks that highly leveraged
entities pose to banks. \term{Archegos operated as a family office rather than a
hedge fund}, but used total return swaps that created exposures similar to those
of a highly leveraged fund. Credit Suisse suffered losses of about \$5 billion,
and other banks collectively lost an additional \$5 billion.
\end{imkeybox}

\begin{imtrapbox}[breakable=false,title={Archegos was not a hedge fund}]
That is the point of including it. The systemic exposure came from the
\emph{structure of the trades} --- total return swaps --- not from the legal form
of the counterparty. A question that calls Archegos a hedge fund has missed why
the case is in the curriculum.
\end{imtrapbox}

% ------------------------------------------------------------------
\lo{IM-9 c}{Summarize different approaches taken by regulators to mitigate risks related to hedge funds.}

\srcnote{The objective the source notes leave unlabelled. The material below is
theirs; only the heading is supplied.}

Hedge funds generally operate with \term{less regulatory oversight than mutual
funds}. As a result, investors are expected to conduct their own due diligence
before investing.

\renewcommand{\arraystretch}{1.2}
\noindent\begin{tabularx}{\textwidth}{@{}>{\raggedright\arraybackslash}p{2.3cm}X@{}}
\toprule
\textbf{\sffamily\small Jurisdiction} & \textbf{\sffamily\small Approach} \\
\midrule
United States & Under the \term{Investment Advisers Act of 1940}, managers were previously exempt from SEC registration if they advised fewer than 15 clients --- the \term{private adviser exemption}. The Global Financial Crisis exposed hedge funds' interconnectedness with banks, and the \term{Dodd-Frank Act (2010)} now requires advisers managing at least \$150 million to register with the SEC, maintain extensive records, appoint a \term{Chief Compliance Officer}, and undergo periodic SEC examinations. \\
European Union & The \term{Alternative Investment Fund Managers Directive (AIFMD)}, 2010, established comprehensive rules: mandatory registration of managers, detailed reporting and disclosure, \term{limits on leverage}, and marketing rules for EU and non-EU funds targeting EU investors. \term{ESMA} was later created to interpret AIFMD provisions, although \emph{enforcement is handled by individual member states}. \\
Singapore & Funds with more than \$250 million in assets must register with the \term{Monetary Authority of Singapore}, provide quarterly unaudited and annual audited reports to investors and MAS, and obtain a \term{Capital Markets Services licence}. \\
Hong Kong & Regulated under the \term{Securities and Futures Ordinance}. Managers must obtain licences depending on their activities and follow reporting and disclosure requirements. \\
China & Funds are classified as \emph{government-supported} or \emph{private}. The \term{CSRC} limits leverage and restricts short selling to certain CSI 300 securities; \term{naked short selling is not allowed} and borrowing shares is expensive. The \term{CIRC} encourages insurers to create asset management firms investing in hedge-fund-like strategies. \\
\bottomrule
\end{tabularx}
\renewcommand{\arraystretch}{1}
\figcap{Five jurisdictions, three different instruments. The United States
regulates the \emph{adviser}; the European Union regulates the \emph{manager and
the leverage}; China regulates the \emph{trade} by restricting which securities
may be shorted at all. Note the split in the EU row: ESMA interprets, member
states enforce --- a classic who-does-what distinction.}

\subsection{Regulatory objectives}

\begin{enumerate}
\item \term{Investor protection.} Ensure hedge funds do not defraud investors.
Participation is typically restricted to \emph{sophisticated} investors capable of
evaluating hedge fund risks.
\item \term{Market integrity.} Prevent illegal practices such as insider trading,
and ensure hedge funds follow the same enforcement standards as other market
participants.
\item \term{Financial stability.} Reduce systemic risk by \term{limiting leverage
and requiring timely posting of collateral}.
\end{enumerate}

\subsection{Mitigating systemic risk, and the cost of overdoing it}

Three approaches are suggested: \term{(1)} more conservative lending by banks;
\term{(2)} lower leverage and greater diversification in hedge funds; and
\term{(3)} reasonable regulation.

\begin{imtrapbox}[breakable=false,title={The qualifier on the third one}]
\term{Overly strict regulation could reduce market liquidity}, because hedge
funds often act as \emph{investors of last resort} in distressed markets. The
curriculum's position is not that more regulation is always better --- it is that
the regulation should be \emph{reasonable}. An option that omits this trade-off
has removed the only judgement in the paragraph.
\end{imtrapbox}

\subsection{Industry self-regulation}

The \term{Alternative Investment Management Association} has contributed by
publishing its \textit{Guide to Sound Practices for Hedge Fund Valuation}.
Publicly listed hedge fund firms such as Man Group and Blackstone disclose
information about their hedge fund activities, giving investors more insight.

% ------------------------------------------------------------------
\lo{IM-9 d}{Describe a typical organizational structure for a hedge fund and explain motivations for adopting this structure.}

Normally, mutual funds must register with the SEC under the \term{Investment
Company Act of 1940}. If registered, they must follow strict rules entailing
regular reporting and limits on leverage, short selling and performance fees.
\term{Hedge funds typically avoid these restrictions by relying on exemptions} ---
and the structure of a hedge fund is largely the shape those exemptions force it
into.

\renewcommand{\arraystretch}{1.2}
\noindent\begin{tabularx}{\textwidth}{@{}>{\raggedright\arraybackslash}p{3.6cm}X@{}}
\toprule
\textbf{\sffamily\small Exemption relied on} & \textbf{\sffamily\small The structural constraint it imposes} \\
\midrule
Section 3(c)(1) & Fewer than \term{100 investors}. \\
Section 3(c)(7) & Only \term{qualified purchasers} --- generally individuals with more than \$5 million in investments. But if such a fund has \term{500 or more investors} it must register under the Securities Exchange Act of 1934. \\
Private adviser exemption (Advisers Act 1940) & Adviser advises \term{fewer than 15 funds}. Many hedge funds structured themselves to remain below this threshold. \\
ERISA avoidance & Pension fund investments kept below \term{25\% of total fund assets}, which prevents the fund from being treated as managing pension assets directly. \\
Regulation D & Capital raised through \term{private placements} rather than public offerings, accepting only \term{accredited investors} --- generally net worth above \$2.5 million or annual income above \$250,000 with expectations of maintaining it. \term{Rule 506} allows up to \term{35 non-accredited investors}. \\
\bottomrule
\end{tabularx}
\renewcommand{\arraystretch}{1}
\figcap{The structure is the exemptions. Every number in the right-hand column is
a threshold a fund is organised to stay beneath, which is why hedge fund
structures look arbitrary until you read them as constraints. Both the United
States and the United Kingdom have historically maintained a relatively favourable
environment, although regulators have moved to control preferential treatment
through \emph{side letters}, which grant certain investors advantages over others.}

\subsection{How the structure changed after the crisis}

\begin{itemize}
\item \textbf{2004:} the SEC tried to force hedge funds to register.
\item \textbf{2006:} the U.S. Court of Appeals for the D.C.\ Circuit
\term{overturned} the rule.
\item \textbf{2007:} the President's Working Group on Financial Markets rejected
strict regulation.
\item Later events --- the \term{Bernie Madoff fraud} and the Global Financial
Crisis --- renewed regulatory pressure, leading to oversight under
\term{Dodd-Frank (2010)}.
\end{itemize}

\begin{imkeybox}[title={What Title IV actually did}]
Title IV of Dodd-Frank \term{eliminated the private adviser exemption}. The
``15-client'' threshold no longer applies and advisers must register with the SEC.
Funds must also report to the \term{Financial Stability Oversight Council} so
regulators can monitor AUM, leverage, counterparty exposures and strategies.

\smallskip
\textbf{Form PF:} funds with at least \$150 million in AUM file \term{annually};
funds with more than \$1.5 billion file \term{quarterly}.
\end{imkeybox}

\subsection{International structures}

Europe follows a stricter approach: funds with more than \texteuro 100 million in AUM
must comply with detailed reporting on leverage, risk exposure and strategies. EU
funds benefit from the \term{EU Passport}, allowing them to market across member
states. Regulations also include disclosure of short positions, restrictions on
naked short selling, and reporting of derivatives transactions. Enforcement varies:
\term{France} focuses on leverage limits, \term{Portugal} closely regulates
derivatives markets, and \term{Germany} restricts bank lending to highly leveraged
investment funds.

\begin{imtrapbox}[title={Consolidated trap summary --- IM-9}]
\begin{itemize}
\item \term{Source.} No GARP book covers this reading. Treat gap-fills as one rung
below a GARP source.
\item \term{Sibling, IM-9 a.} Portfolio-level risks are benchmarking,
survivorship bias and UBTI. Everything else in the list --- counterparty, key
person, style drift, correlation spikes --- is investment-level.
\item \term{Role, IM-9 b.} Revenue risk needs no default at all; counterparty risk
needs one. Credit risk sits between them and turns on collateral value.
\item \term{Definition, IM-9 b.} \term{Archegos was a family office}, not a hedge
fund. The exposure came from total return swaps.
\item \term{Role, IM-9 b.} The Bank of England's Amaranth verdict was that
\emph{major intermediaries contribute more systemic risk than hedge funds do}.
\item \term{Role, IM-9 c.} \term{ESMA interprets AIFMD; member states enforce it.}
\item \term{Scope, IM-9 c.} Overly strict regulation \emph{reduces} market
liquidity, because hedge funds act as investors of last resort.
\item \term{Number, IM-9 d.} 100 investors for 3(c)(1); 500 triggers 1934 Act
registration for a 3(c)(7) fund; 25\% ERISA threshold; 35 non-accredited investors
under Rule 506; \$150m annual and \$1.5bn quarterly for Form PF. These are the
most corruptible facts in the reading.
\end{itemize}
\end{imtrapbox}

% ==================================================================
